\section{Einleitung}\label{sec:Einleitung}

\subsection{Motivation}\label{sec:Motivation}
Durch stetige Innovationen und eine extreme Globalisierung, hat sich die Wirtschaft stark
verändert. Der moderne Markt erfordert ein schnelles Reagieren auf Trends und Veränderungen.
Dies führt zu einem massiven Forschungs- und Innovationsdruck, der sowohl Vorteile als
auch Nachteile mit sich zieht. 
Aus diesem Grund entwickeln sich heutzutage viele Bereiche der Industrie hin zu einer effizienteren und
schnelleren Produktion. Dies erfordert eine stetige Weiterentwicklung sowohl der
Fertigungsverfahren als auch der Produkte selbst.\\
Dieser Prozess lässt sich besonders mit der Transformation im Mobilitätssektor aufzeigen.
Die Elektrifizierung des Verkehrs entstand ursprünglich aus einer Vision des umweltbewussteren
Transports. Dieser Gedanke wurde zunächst in Forschungsprojekten verfolgt, um eine Machbarkeit
von E-Fahrzeugen zu demonstieren. Parralel dazu gab es massive Fortschirtte in der Batterietechnik,
Leistungselektronik, Antriebstechnik und Materialforschung, die alle dazu betragen sollten, dass 
Konzept des E-Autos Alltagstauglich ist.\\
Hier begann nun eine Entwicklung, die sehr Typisch für moderen Wirtschaftstrends sind: Die
etablierten Anbieter von Fahrzeugen mit Verbrennungsmotren blieben bei ihrer Technologie und sind
kein Risiko eingegangen. Auf der anderne Seite haben neue Firmen, besonders aus China, die
Chance gesehen im globalen Markt mit anzutreten. Mittlerweile hat sich diese Dynamik dahin entwickelt,
das es einen massiven Kampf jeglicher Unternehmen, um Marktanteile in der ganzen Welt, gibt.\\
Aktuelle Zahlen zeigen, dass auch in Deutschland die Zulassung von Eletkrisch angetriebenen Verkehrsmittel
stetig zunimmt. Dies hat nicht nur einen ökologischen Hintergrund, sondern auch wirtschaftliche und praktische
Gründe. So ist die Ladeinfrastruktur und Ladezeit in einem alltagstauglichen Zustand, aber
auch Anschaffungs- und Laufende Kosten sind durch Förderungen und einem massiven Preisdruck durch
die Hersteller, vergleichbar mit Autos mit konventionellen Antrieben. Dieser Preisdruck führte zu
einer intensiven und schnellen Weiterentwicklung. Hierbei wurden nicht nur die Subsysteme optimiert, 
sondern auch die Art der Produktion.\\ \todo{Übergang?, was ist mit Automatisierung}
Der große Bedarf an solchen Systemen zeigt sich z.B. der Entwicklung von modernen elektrischen Antrieben,
die genau diesem Muster folgen. Hierbei spielt die Hairpin-Technologie eine entscheidende Rolle, da sowohl die
Produktion als auch das Endprodukt derzeit stark optimiert werden. Im Gegensatz zur
Runddraht-Wickeltechnologie werden bei der Hairpin-Technologie rechteckig vorgeformte
Kupferleitungen statt gewickelter Spulen im Stator eingesetzt. Dieser Prozess bietet Vorteile,
wie eine erhöhte Kupferdichte, bessere Kühlung, und eine vereinfachte Montage. In der
Prototypenproduktion werden zur Herstellung von Hairpins Kupferleitungen präzise gebogen
und abisoliert. Da die derzeitige Entwicklung eine hohe Stückzahl verlangt, ist eine
weitgehende Automatisierung hilfreich. Insbesondere die Ablösung des manuellen
Abisolierens in der Prototypenfertigung bietet das Potenzial, Geschwindigkeit, Kosten und
Sicherheit in der Entwicklung zu optimieren.
%
\subsection{Zielsetzung}\label{sec:Unterkapitel}
Was soll erreicht werden?
%
\subsection{Aufgabenstellung}\label{sec:Unterkapitel}
Es muss ein System entwickelt werden, welches die Hairpinenden automatisiert und
präzise abisoliert. Das Gesamtsystem besteht aus einer Laseranlage zur Entfernung der
Isolierung der Hairpins, einem Roboterarm zur Handhabung der Hairpins und einer
Zuführung. Der Fokus dieser Arbeit liegt auf der Konzeption, Realisierung und dem Testen
des Zuführsystems.
Um diese Ziele zu erfüllen, wird in mehrere Arbeitspakete unterteilt. Zunächst muss sich in
bestehende Systeme eingearbeitet und ein Konzept erstellt werden. Weiterführend muss
dieses Konzept konkretisiert und ausgearbeitet werden. In einem realen System wird das
entwickelte Konzept aufgebaut. Als letztes folgt die Test- und Optimierungsphase. Daraus
leitet sich also die Forschungsfrage ab: Wie kann ein System zur automatisierten
Abisolierung von Hairpins in der Prototypenproduktion aussehen?
Was muss gemacht werden?
%

\subsection{Aufbau der Arbeit}\label{sec:Unterkapitel}
Wie wird die Arbeit erledigt?
%